% OGC2026 Technical Report.  Anonymous: the contributions section uses "Member A/B" by design.
\documentclass[runningheads]{llncs}

\usepackage[T1]{fontenc}
\usepackage{graphicx}
\usepackage{amsmath}
\usepackage{booktabs}
\usepackage[protrusion=true,expansion=false]{microtype}

\begin{document}

\title{OGC2026 Technical Report}
\maketitle

\begin{abstract}
The challenge asks for a schedule and a physical layout at the same time: every block needs a
bay, an entry time, and a pose that a vertically descending crane can actually reach. Our
starting observation is that the objective $w_1Z_1+w_2Z_2+w_3Z_3$ reads only the pair (bay,
entry time)---coordinates and orientation appear nowhere in it. Geometry is therefore a
feasibility constraint that decides how much scheduling freedom exists, not a term to be optimised
directly. The two training sets do not even share a shape --- over half the preliminary instances
carry no tardiness at all, none of the final forty do --- so we avoid density thresholds and instance
classification entirely: one portfolio of six operators competes for the budget on
measured payoff per second, and every hot loop lives in one of two C++ extensions. Our main contributions are an exact contact upper bound that skips $40.8\%$ of the
candidate cells while provably returning the identical solution, a $10.5\times$ faster
conflict-graph build verified by an unchanged edge count, and a position-independent geometric
predicate that removed a floating-point inconsistency from the packing test. The most instructive
result is negative: that same verified $1.85\times$ speedup made one instance $7.7\%$ worse,
because the operator schedule runs on a wall clock and a faster search therefore takes a different
path rather than a longer one.
\end{abstract}

\section{Introduction}

\subsection{Challenge Problem \& Main Difficulties}

Assigning hull blocks to bays over time is the spatial scheduling problem of shipbuilding. Three
properties of the version posed here shaped everything we built.

\paragraph{The objective and the constraints read different variables.}
Tardiness, workload imbalance and preference are all functions of which bay a block occupies and
when it enters; $x$, $y$ and orientation appear in none of them. Geometry enters only through
feasibility. A layout that cannot accommodate a block pushes its entry later, and that is the only
channel through which packing quality becomes objective value. Treating tightness as an objective
term is a modelling error we made early. Weighting contact more heavily does not buy a better
schedule, it buys a denser layout whose scheduling consequences are unmeasured.

\paragraph{The crane makes the packing three-dimensional.}
A block is lowered vertically, so a placement is legal only if every layer of the descending block
clears everything present along the whole descent. Two neighbours occupying the same floor cells
but standing one and two layers high ration space completely differently. The scarce resource is
height, not occupancy, and that is why the descent test goes layer by layer.

\paragraph{The two rounds do not have the same shape.}
At a uniform $60$\,s the preliminary training set is bimodal: $Z_1$ is exactly zero on twenty-one of
its forty instances and strictly positive on the rest (Table~\ref{tab:train}). The final-round
set is not bimodal at all. All forty of its instances are late, and the share of the objective
carried by tardiness spreads from $2.5\%$ to $97.9\%$ with no gap wider than fifteen points
anywhere in it (Table~\ref{tab:final}). Where $Z_1$ vanishes the entry times are free and the problem is an
assignment dominated by preference; where the yard is saturated they are the whole objective.
Neither shape admits a threshold --- in the first the classes are separate but the boundary is
unobservable, in the second there are no classes --- so we built a scheduler that measures which
operator is paying rather than one that classifies the instance.

\subsection{Contributions}

\begin{itemize}
\itemsep0pt
\item A \emph{gate-free operator portfolio}: six operators compete for the budget on realised
objective gain per second, with no density threshold anywhere; slice sizing is a separate question
answered by completion.
\item An \emph{exact contact upper bound} skipping $40.8\%$ of candidate cells while provably
returning the identical solution, verified per cell rather than by comparing run outcomes.
\item A $10.5\times$ faster \emph{conflict-graph build}, accepted on an unchanged edge count rather
than a stopwatch, over a shared memo made safe by compare-and-exchange.
\item \emph{Offset-relative geometry}: the packing predicate became position-independent,
fixing a disagreement we had misdiagnosed as hashing.
\item A negative result more useful than any of these: an exact, verified $1.85\times$ speedup that
made one instance $7.7\%$ worse (Section~\ref{sec:wallclock}).
\end{itemize}

\section{Algorithm Description}

\subsection{Overall Algorithm Framework}

Four stages, and no preprocessing worth the name. \emph{Initial solution:} each of four independent
worker processes seeds a pool with a cheap always-feasible sequential schedule --- the one guarantee
that the algorithm never returns nothing, and the reason every other fallback path could be
deleted. \emph{Improvement:} each worker then runs an operator loop until its budget is exhausted,
differing from the others in random seed and in its starting offset into six diversification axes,
so the four explore different trajectories over the same operators. \emph{Selection:} the best full
objective across workers is taken, and the entry point reserves $\min(0.2T, 40\,\mathrm{s})$ of the
limit for a final preference-repair pass over that winner. \emph{Validation:} feasibility is
enforced inside the operators and re-checked outside them --- one acceptance function re-runs the
provided checker on anything that beats the incumbent and returns $+\infty$ if it fails. An operator
may be wrong about feasibility; a wrong answer cannot be accepted.

There are six operators. A contact-beam construction. An evolutionary re-growth with crossover and
path-relinking over (bay, dispatch-order) anchors, decoded by the same beam. A workload-balance
repair, a preference repair, a CP-SAT bay reassignment, and the exact bay repacking of
Section~\ref{sec:repack}. Each receives one probe to establish a rate, and the budget then goes to
whichever has the highest realised objective gain per second. Fifteen percent is held back for
exploration so a slow starter can recover. The arrangement is an adaptive large neighbourhood
search~\cite{ropke2006}, with one difference. Its adaptive weights score realised objective gain
per second, where the usual form counts accepted moves.

\paragraph{Selection and sizing are different questions.}
Slice size is deliberately not driven by the selection signal, because sizing on payoff spirals
downwards. A construction that returns a good solution but merely fails to beat the incumbent
scores zero gain. It shrinks, and then it can no longer finish at all. We measured fourteen calls returning four
solutions while the objective moved from $48{,}840$ to $56{,}841$. Size is a completion question instead: an operator
that returned nothing was starved and is given more, one that returned anything fitted, and a
repair pass is given what it used.

\subsection{Key Ideas that Succeeded}
\label{sec:succeeded}

Five changes moved the objective. They are given in the order they act, from the construction that
produces a solution to the operator that repacks one.

\paragraph{Contact beam, and why its contact measure is flat.}
The construction is a filtered beam search over dispatch order~\cite{ow1988}, ranking survivors on
the exact objective rather than a priority index. Each state places the next block at the position
minimising a weighted sum of negative contact, a positional sweep term and a preference penalty,
and survivors are completed by a rollout. Contact plays the role touching perimeter plays in
two-dimensional packing, extended to a descent the crane must clear.

Contact admits two measures. A per-layer one describes the geometry better --- a block beside a
shorter neighbour earns nothing for the layers facing open air --- and it is implemented, but the
submitted configuration uses the flat measure because the bound below is derived against a flat
occupancy prefix sum, which does not dominate per-layer occupancy. The two cannot both be on, and
the bound is worth $1.09$--$1.85\times$. This is a constraint we accepted rather than a property we
wanted; with a per-layer bound the choice would likely go the other way.

\paragraph{An exact contact upper bound.}
Profiling showed $86\%$ of scored cells were evaluated after the best cell of that scan had already
been found. Bounding the positional term fired zero times in $112.9$ million cells. The score is
contact-dominated by more than an order of magnitude, with ranges of about $2.3$ against about
$60$, so a bound that knows only the sweep coordinate is never competitive.

Bounding the contact itself inverts this. Contact counts the four neighbours of each occupied
footprint cell. A neighbour outside the bay scores one, a neighbour inside the same footprint
scores nothing, and an occupied neighbour scores its weight. Every cell that can contribute
therefore lies inside the footprint bounding box dilated by one, and counts at most once per
direction, giving
\begin{equation}
\mathrm{ct}(i_x,i_y) \;\le\; W(i_x,i_y) \;+\; 4\,w_{\max}\,
\bigl|\,\mathcal{O} \cap B^{+1}(i_x,i_y)\,\bigr| ,
\end{equation}
where $W$ counts cell--direction pairs pointing out of the bay, $\mathcal{O}$ is the occupied set
and $B^{+1}$ the dilated box. An occupancy prefix sum, built once per (bay, time window) at the
same asymptotic cost as the occupancy map it accompanies, answers the rectangle in four reads, and
a cell whose best possible score cannot beat the incumbent is skipped without any geometric test.
The bound is applied only where it is proved, which is what excludes the per-layer variant.

\paragraph{A cliff at the bottom of the beam, and the axis that came out of it.}
The beam checked its deadline at each level and returned nothing when it passed, so every worker
fell back to the sequential floor. The failure was not gradual. On a $300$-block instance the
answer was $4.02\times10^9$ at $110$\,s against $9.69\times10^7$ at $130$\,s, a factor of $42$, and
$143$ on a $250$-block one, at a threshold that moves with instance size while the hidden limits
are undisclosed. The repair needed no new machinery, since the beam already completes a partial
state by rollout to rank survivors: on expiry it completes the incumbent rather than discarding
it.

That exposes a knob rather than a patch: the fraction of its slice the beam claims before handing
the remainder to the rollout. Sweeping it down improves large instances monotonically ($-24.3\%$ to
$-11.9\%$ at one tenth) and harms small ones monotonically, since the beam finishes anyway there
and a low aim only narrows it ($+25.09\%$ at $150$ blocks). We did not resolve that by testing
instance size. The workers are already a portfolio whose minimum is the answer, so half run the low
aim and half the high one, each group carrying the instances the other loses, and nothing is gated.
Over $37$ paired final instances this is $31$ better against $6$ worse, median $-3.44\%$, sign test
$p<0.0001$, splitting $17$--$2$ on the large group and $14$--$4$ on the small. Our prediction that
the small group would pay was wrong: true of a constant, false of an axis. A portfolio does not
need every member to be good, it needs them to fail on different instances.

Replicating the two small instances that lost by more than $10\%$ shrank both to $+3.2\%$. A single
paired draw at $150$ blocks is worth its sign, not its magnitude.

\paragraph{A conflict-graph build that watches its own clock.}
The repacking operator --- which lifts every block out of one bay and re-seats the whole set
exactly, and is described in Section~\ref{sec:repack} --- begins with an $O(n_{\text{col}}^2)$
conflict graph that cannot be interrupted. Two columns conflict only while both are present, so sorting by entry time and
breaking the inner loop once the later entry passes the earlier exit skips $83\%$ of pairs without
visiting them, and the fields the filter reads moved into flat arrays. Neither change can alter
which pairs conflict. The edge count therefore had to come out identical, and that was the
acceptance test, not the clock. The same configuration now builds in $10.7$--$14.2$\,s against about $90$\,s
before. The build also projects its own completion every $256$ rows and steps down to a cheaper
configuration it can finish rather than losing the call.

The rows are independent, so the loop is finally parallel. Where the results go is not independent.
Each thread keeps its own edge list, and the memo on the conflict predicate is shared.
Sharing it naively is not merely racy but silently wrong, because the table is open-addressed and
two threads inserting different keys can both claim one empty slot, the loser then writing its
verdict into the winner's entry. A slot is therefore claimed by compare-and-exchange, with relaxed
ordering, since the value is a pure function of the key. A private memo per thread gave almost all
the parallel gain back ($1.10\times$), reuse being where this loop's speed lives, while sharing it
correctly reaches $1.43\times$ on four cores at an edge count of $133{,}863{,}098$ in every arm.

An identical edge count is also what made the memo usable: its first version produced $36$ extra
edges, and dumping the disagreements showed the predicate itself was position-dependent through
floating-point rounding, not the key aliasing we assumed. Comparing origin outlines plus an offset
makes the verdict a function of relative position alone, fixing the memo by fixing the predicate.


\paragraph{Pricing a seat, not a block.}
The operator hands the solver a list of candidate placements, a \emph{column} being one block at
one orientation, position and entry time, and asks for the maximum-weight conflict-free selection.
The weight was attached to the \emph{block}. Since the objective reads only (bay, entry time) and
one call concerns one bay, the entry variant is exactly the granularity at which a block's
alternatives differ: an on-time column and a ten-days-late column into a preferred bay carried the
same number, so the solver took whichever seated more easily.

Every attempt to express the trade from outside failed for the same reason. Discounting a block's
weight by the cost of its late seat undervalues it when it takes the on-time one, and at the
break-even delay, the only principled place to put the offer, the discount equals the entire gain:
the weights of precisely the high-regret blocks collapse to the floor. Across the final instances,
discounting is indistinguishable from never offering a late window at all: $9$ better against $21$
worse, median $+0.00\%$.

Attaching the weight to the entry variant removes the workaround instead of tuning it. Each offered
window is priced by the only term an entry time can move,
\begin{equation}
p(w) \;=\; g \;-\; w_1\bigl(\max(0,\,e_w - d) - \max(0,\,e_0 - d)\bigr),
\end{equation}
with $g$ the gain of the move, $e_w$ the exit under window $w$, $e_0$ the current exit and $d$ the
due date. The expression is exact in both directions. A window \emph{earlier} than the block's
current entry therefore prices above the base, and the solver can be paid for pulling a block
forward. A per-block weight cannot represent that at all. On $39$ final instances the priced form is
$20$ better against $9$ worse (sign test $p\approx0.03$, mean $-0.74\%$), and it does no harm where
no late window can open.

\subsection{Key Ideas that Failed}
\label{sec:failed}

\paragraph{A wall-clock schedule turns a speedup into a different answer.}
\label{sec:wallclock}
This is the result we would most like to pass on. Having verified that the contact upper bound is
exact, that it returns bit-identical placements at fixed work, and that it is $1.85\times$,
$1.66\times$ and $1.09\times$ faster on three instances of very different density, we enabled it
and one instance got $7.7\%$ worse.

The bound is not at fault, and neither is anything else in the search. The operator loop decides
what to run next, and for how long, from elapsed seconds, so a faster scan fits more calls into the
same budget and shifts the diversification axis rotation, the solution pool, and the repacking
operator's call counter and hence its seed. That operator is handed a different incumbent, and on
this instance exactly one incumbent, worth $107{,}195$, ever leads it to $83{,}095$, which is in
turn the only route to the best solution we have recorded there. Every run reaching that solution
had $83{,}095$ in its log and every run that did not, did not. With the bound off we reproduce it
three times in four; with it on, never in two.

The lesson generalises beyond this bound. Making a time-budgeted portfolio faster does not make it
search longer along the same path, it makes it take a different path, and the new one need not be
better. We priced that per instance, as the competition scores, and shipped with the bound
disabled: $-7.7\%$ against $+2.2\%$ is a losing trade.

\paragraph{A better incumbent can repack worse.}
Incumbents worth $97{,}570$, $100{,}685$ and $107{,}195$ yielded repacking improvements of
$9{,}580$, $13{,}510$ and $24{,}100$. A tight solution has no slack left, and this operator eats
slack. Its result follows the incumbent's layout, not its objective. The roster nevertheless hands
it the best-scoring solution every time.

\paragraph{Three changes that did not pay.}
The headroom from the faster conflict-graph build bought nothing: widening the $40$-outsider
candidate cap became affordable ($331$\,s of a $180$\,s budget before, $179$\,s after) and scored
$3$ better against $4$ worse, median $+0.17\%$ --- the evidence for wanting it had come from two
instances.
A free-column bitset for the packing search gave identical results and measured
$1.6$--$1.9\times$ \emph{slower}: its saving scales with columns per block, its cost with edges, and
we had assumed the wrong line was hot. We also twice predicted the contact bound's benefit from its
firing rate and were wrong in the same direction both times --- most cells are rejected by a cheap
bitmap probe, while the cells a bound removes are exactly those that would have reached the exact
geometric test. Cell counts are not seconds.

\subsection{Further Improvement Plan}
\label{sec:plan}

The first item is to remove the wall clock from the schedule's decisions. We have localised where
it enters. With the deadline lifted, the construction returns identical placements with and without
the bound; under the deadline it receives, the two diverge. The divergence therefore begins in the
construction's stopping rule, not in the operator loop. The packing solver already takes an
iteration cap for this reason. Once the schedule depends on work done instead of seconds elapsed,
the contact bound can be re-enabled and contribute its speed rather than a trajectory.

Second, the repacking operator should sample the solution pool instead of always taking its best
member. Its yield depends on slack, not on quality. The re-growth operator already samples the
pool, so the mechanism is there to reuse.

Third, that operator's cost model deserves the scrutiny its search has had. It declines a
configuration whose predicted build will not fit the budget, and the prediction charges a fixed
parsing cost to a quadratic term and mixes two different column counts. Both errors are squared,
and both push towards declining work the budget could afford.

\section{Implementation Details}
\label{sec:implementation}

Python~3.12 orchestrates and two C++17 extensions built with \texttt{pybind11} and OpenMP hold
every hot loop. CP-SAT serves one operator and the algorithm degrades to the remaining five if it
is unavailable. Python holds only control flow --- the scheduler, the operator definitions, the
acceptance test --- while every geometric predicate, every candidate scan and both search engines
are C++, the construction parallelising over beam states and the portfolio over worker processes.

\paragraph{The repacking operator.}
\label{sec:repack}
This operator lifts every block out of the most contested bay, adds the outside blocks that might
want to enter it, and hands the set to a maximum-weight set-packing solver over space--time columns
weighted in objective units. Destroying a bay and rebuilding it is ruin and
recreate~\cite{schrimpf2000} with an exact recreate step, and the selection is a maximum-weight
independent set searched by the swap-based scheme of Andrade et al.~\cite{andrade2012}. A call has
two phases of different character: an uninterruptible build and a deadline-honouring search. The
solver takes a total-time bound and enforces it against its own measured build cost rather than a
predicted one; an earlier version subtracted a predicted build from the search budget as well,
double-charging it. The configuration ladder is expressed as fractions of each instance's own
entry-time ceiling, so window widths come from the instance rather than being fitted to it.

\paragraph{Verification instrumentation.}
The engines carry counters for cells visited, exact tests, cheap rejections and cells skipped by
the bound, plus a soundness mode that scores every cell the bound wanted to skip and counts those
that would have won: zero out of $45.9$, $11.6$ and $9.3$ million. An earlier bound had passed an
identity check that was meaningless because the branch never executed, so we stopped accepting
outcome comparisons as evidence.

\section{Computational Results}

Table~\ref{tab:train} reports all forty preliminary-round training instances at a uniform $60$\,s
budget, decomposed into the three objective components; a uniform budget is the only comparable
one, since the hidden limits vary and are not disclosed. All forty pass the provided feasibility
checker, and the split is visible directly: $Z_1$ is zero on every instance from 1 to 20 and
positive on the rest bar instance 29, where this build also reaches zero.
Table~\ref{tab:final} puts the two rounds side by side.

\begin{table}[t]
\centering
\caption{Preliminary-round training set, all forty instances at $60$\,s. Sizes and bay counts are
in Table~\ref{tab:final}.}
\label{tab:train}
\footnotesize
\setlength{\tabcolsep}{1pt}
\begin{tabular}{rrrrr@{\;}rrrrr@{\;}rrrrr}
\toprule
\# & $Z_1$ & $Z_2$ & $Z_3$ & Obj. & \# & $Z_1$ & $Z_2$ & $Z_3$ & Obj. & \# & $Z_1$ & $Z_2$ & $Z_3$ & Obj. \\
\midrule
1 & 0 & 157 & 2 & 1{,}499 & 15 & 0 & 4{,}366 & 155 & 42{,}445 & 29 & 0 & 6{,}364 & 1{,}212 & 382{,}692 \\
2 & 0 & 144 & 15 & 3{,}690 & 16 & 0 & 3{,}769 & 145 & 38{,}130 & 30 & 123 & 573 & 1{,}442 & 1{,}930{,}651 \\
3 & 0 & 2{,}281 & 137 & 43{,}360 & 17 & 0 & 6{,}575 & 213 & 54{,}629 & 31 & 233 & 2{,}807 & 10{,}151 & 5{,}825{,}327 \\
4 & 0 & 2{,}278 & 0 & 15{,}946 & 18 & 0 & 4{,}098 & 218 & 45{,}386 & 32 & 411 & 1{,}925 & 2{,}337 & 2{,}781{,}688 \\
5 & 0 & 1{,}834 & 228 & 47{,}038 & 19 & 0 & 4{,}350 & 321 & 60{,}093 & 33 & 839 & 1{,}666 & 4{,}591 & 6{,}298{,}923 \\
6 & 0 & 2{,}746 & 133 & 39{,}172 & 20 & 0 & 3{,}693 & 618 & 99{,}408 & 34 & 102 & 1{,}970 & 1{,}368 & 1{,}082{,}900 \\
7 & 0 & 2{,}514 & 254 & 55{,}698 & 21 & 19 & 1{,}530 & 2{,}202 & 598{,}927 & 35 & 26 & 3{,}991 & 1{,}547 & 598{,}663 \\
8 & 0 & 2{,}413 & 8 & 11{,}252 & 22 & 1 & 2{,}955 & 1{,}792 & 738{,}998 & 36 & 28 & 4{,}383 & 6{,}813 & 111{,}628 \\
9 & 0 & 2{,}147 & 265 & 50{,}485 & 23 & 90 & 93 & 1{,}560 & 1{,}532{,}961 & 37 & 585 & 1{,}783 & 3{,}579 & 4{,}104{,}337 \\
10 & 0 & 1{,}419 & 450 & 69{,}783 & 24 & 2 & 4{,}235 & 717 & 262{,}941 & 38 & 2{,}306 & 3{,}683 & 5{,}788 & 32{,}489{,}664 \\
11 & 0 & 305 & 137 & 20{,}356 & 25 & 266 & 2{,}772 & 1{,}662 & 213{,}434 & 39 & 391 & 1{,}219 & 5{,}062 & 5{,}977{,}379 \\
12 & 0 & 2{,}385 & 450 & 76{,}545 & 26 & 440 & 1{,}216 & 5{,}100 & 6{,}640{,}032 & 40 & 2{,}216 & 3{,}156 & 8{,}947 & 1{,}597{,}539 \\
13 & 0 & 3{,}327 & 476 & 79{,}943 & 27 & 1{,}442 & 22 & 4{,}182 & 20{,}899{,}030 &  &  &  &  &  \\
14 & 0 & 1{,}605 & 280 & 45{,}265 & 28 & 57 & 1{,}385 & 1{,}963 & 1{,}353{,}036 &  &  &  &  &  \\
\bottomrule
\end{tabular}
\end{table}

\begin{table}[t]
\centering
\caption{The two training sets compared, both at $60$\,s on the same build. They are different
instances: $w_1$ spans $333$ to $13{,}559$ on the final set against $667$ to $29{,}630$ on the
preliminary one, and $w_3$ reaches $800$ against $600$.}
\label{tab:final}
\footnotesize
\begin{tabular}{lrrrr}
\toprule
 & \multicolumn{2}{c}{Preliminary} & \multicolumn{2}{c}{Final} \\
\cmidrule(lr){2-3}\cmidrule(lr){4-5}
 & median & range & median & range \\
\midrule
Blocks $n$ & 200 & 100--300 & 225 & 150--300 \\
Bays $m$ & 3 & 2--5 & 3 & 2--5 \\
Tardiness $Z_1$ & 0 & 0--2{,}306 & 480 & 1--10{,}536 \\
Imbalance $Z_2$ & 2{,}333 & 22--6{,}575 & 2{,}425 & 91--15{,}458 \\
Preference $Z_3$ & 668 & 0--10{,}151 & 4{,}210 & 196--10{,}614 \\
Objective & 105{,}518 & 1{,}499--32{,}489{,}664 & 5{,}005{,}623 & 6{,}804--75{,}792{,}403 \\
\midrule
Instances with $Z_1=0$ & \multicolumn{2}{c}{21 of 40} & \multicolumn{2}{c}{0 of 40} \\
Feasible & \multicolumn{2}{c}{40 of 40} & \multicolumn{2}{c}{40 of 40} \\
\bottomrule
\end{tabular}
\end{table}

\begin{table}[t]
\centering
\caption{Objective returned on the eight final-round instances by each of our three submissions.
Between the first and the second we shipped per-seat window pricing and the parallel
conflict-graph build; between the second and the third, the beam salvage of
Section~\ref{sec:succeeded} and the split-aim worker portfolio.}
\label{tab:subs}
\begin{tabular}{lrrrr}
\toprule
Instance & Submission 1 & Submission 2 & Submission 3 & 1 to 3 \\
\midrule
1     & 3{,}068{,}862  & 3{,}328{,}237  & 2{,}847{,}060  & $-7.2\%$ \\
2     & 19{,}829{,}534 & 20{,}337{,}489 & 20{,}063{,}787 & $+1.2\%$ \\
3     & 5{,}886{,}589  & 5{,}867{,}652  & 5{,}613{,}271  & $-4.6\%$ \\
4     & 4{,}421{,}375  & 4{,}283{,}430  & 4{,}681{,}440  & $+5.9\%$ \\
5     & 6{,}308{,}099  & 6{,}104{,}015  & 6{,}148{,}480  & $-2.5\%$ \\
6     & 1{,}024{,}046  & 1{,}053{,}770  & 968{,}674      & $-5.4\%$ \\
7     & 16{,}385{,}030 & 16{,}285{,}457 & 16{,}310{,}765 & $-0.5\%$ \\
8     & 17{,}079{,}881 & 16{,}023{,}014 & 17{,}347{,}930 & $+1.6\%$ \\
\midrule
Total & 74{,}003{,}416 & 73{,}283{,}064 & 73{,}981{,}407 & $-0.0\%$ \\
\bottomrule
\end{tabular}
\end{table}

\paragraph{Across submissions.}
Table~\ref{tab:subs} gives what each of our three submissions returned on the eight final-round
instances. It does not separate the changes from the noise. Every per-instance move is smaller than
the spread the algorithm produces against itself; three replicates of one unchanged build spread
$29.3\%$, against a largest move here of $7.2\%$. Both shipping steps also bundled two changes.
What the table does record is direction. Five of eight instances improved from the first submission
to the third, and the total is flat.

\paragraph{Per-seat pricing, on the final practice set.}
The change described in Section~\ref{sec:succeeded} was judged on all forty final practice
instances, paired, at $180$\,s, against two controls: no late window offered at all, and the late
window offered under the old per-block weight. Judging by sign count over the set rather than
instance by instance is deliberate. A single instance answers the same question with a spread of
about $3\%$, and unchanged code shifted $8.4\%$ between two runs an hour apart on this machine, so
no single pair can resolve an effect of this size.


On the $35$ instances whose weights permit a late window at all, the previous per-block weight
scored $9$ better against $21$ worse (mean $+0.62\%$) and the per-seat price $20$ against $9$
(mean $-0.74\%$); on the four where no window can open it is $2$--$1$ and does no harm, which was
the condition we had set in advance for enabling it by default. The effect shrank as the sample
grew, from $-1.61\%$ at twelve instances to $-0.74\%$ at thirty-nine: the first twelve had been
chosen as the most favourable.

Two cautions belong with these numbers rather than after them. Run-to-run spread on the densest
instance is about $112{,}000$ on unchanged code, so smaller differences are noise. The sparsest is
not stable at all: its objective is a minimum over the repacking operator's attempts within a run,
that operator's improvement varies from $6{,}590$ to $31{,}315$ with the incumbent it is handed, and
the value we quote appeared in four runs of seven.

\paragraph{Speed, measured at identical work.}
A fixed-time A/B cannot evaluate an accelerated search: both arms consume the same budget, so the
faster one legitimately reaches a different solution. We replay one construction call with
identical arguments and no deadline, comparing a digest over every returned placement. On a sparse,
a mid-density and a dense instance the bound skipped $40.8\%$, $3.4\%$ and $1.3\%$ of cells, none
unsound, for $1.85\times$, $1.66\times$ and $1.09\times$ at an identical digest.


The conflict-graph build was accepted the same way, on an edge count confirmed unchanged on every
configuration tested, for $10.5\times$ in isolation.

\section{Team Members' Contributions}

The rules require anonymity, so members are identified by role. This section states who owned which
decisions; Section~\ref{sec:ai} states how they were carried out. \textbf{Member A} owned the
problem analysis and the algorithmic direction: treating geometry as a constraint on the objective
rather than a term in it, choosing a measured-payoff portfolio over instance classification, setting
what the contact bound had to satisfy before it could ship, and the reading of
Section~\ref{sec:wallclock}. \textbf{Member B} owned the codebase and the experimental
discipline: version management, packaging, the rule that every comparison runs on an unloaded
machine, which measurements were trustworthy enough to act on, and the check that every reported
number matches the log behind it. Every change that shipped, and every change that did not, was
decided by the authors.

\section{AI Use Declaration}
\label{sec:ai}

One generative AI tool was used: Anthropic Claude, through its command-line coding interface, as a
coding and analysis assistant throughout development. The authors set the direction and
the tool did the construction under it. Working from specifications settled in discussion, it
implemented the two C++ extensions and the Python scheduler, worked out the algebra of the contact
bound and the conditions under which it holds, built the measurement harnesses, ran the sweeps,
diagnosed defects among them the position-dependent geometric predicate, and drafted this report
for the authors to revise.

Nothing was adopted on the tool's word. Every change was accepted or rejected by the authors against
a measurement specified in advance, and several were rejected. The numbers here come from executing
the submitted code, not from a language model.

\begin{thebibliography}{4}

\bibitem{andrade2012}
Andrade, D.V., Resende, M.G.C., Werneck, R.F.: Fast local search for the maximum independent
set problem. J. Heuristics \textbf{18}(4) (2012)

\bibitem{ow1988}
Ow, P.S., Morton, T.E.: Filtered beam search in scheduling. Int. J. Prod. Res. \textbf{26}(1) (1988)

\bibitem{ropke2006}
Ropke, S., Pisinger, D.: An adaptive large neighborhood search heuristic for the pickup and
delivery problem with time windows. Transp. Sci. \textbf{40}(4) (2006)

\bibitem{schrimpf2000}
Schrimpf, G., Schneider, J., Stamm-Wilbrandt, H., Dueck, G.: Record breaking optimization
results using the ruin and recreate principle. J. Comput. Phys. \textbf{159}(2) (2000)

\end{thebibliography}

\end{document}
